\section{Experimental Setup}

We design experiments to answer five research questions corresponding to the five validated sub-hypotheses:

\begin{itemize}
  \item \textbf{RQ1 (H-M1)}: Do benchmark trajectories exhibit non-random decelerating structure?
  \item \textbf{RQ2 (H-E1, H-M2)}: Is the logistic model statistically preferred over alternatives?
  \item \textbf{RQ3 (H-M3)}: Are fitted parameters physically plausible?
  \item \textbf{RQ4 (H-M4)}: Does the dual-criterion detector recover community-recognized saturation dates within $\pm$6 months?
  \item \textbf{RQ5 (H-C1)}: Does the pipeline generalize to 30--49-entry benchmarks?
\end{itemize}

\subsection{Benchmarks}

\begin{table}[h]
\centering
\caption{Benchmark dataset summary.}
\begin{tabular}{lllll}
\toprule
Benchmark & Release & Entries & Ground Truth Date & Source \\
\midrule
GLUE & Feb 2019 & 62 & Sept 2019 & \citet{wang2019superglue} \\
SuperGLUE & May 2019 & 87 & Jun 2021 & \citet{srivastava2022bigbench} \\
\bottomrule
\end{tabular}
\end{table}

Data: curated historical scores from published model papers (BERT, XLNet, RoBERTa, T5, DeBERTa, etc.).
Papers With Code API deprecated; curated fallback used (see limitations).

\subsection{Baselines for RQ2}

\textbf{Linear} ($y = at + b$): null hypothesis of constant-rate improvement.
\textbf{Power-law} ($y = at^b$): sub-linear alternative with no asymptote.
Both are standard growth models; their inclusion makes the AIC comparison a genuine three-way test.

\subsection{Synthetic Benchmarks for RQ5}

20 synthetic timeseries at $n \in \{30, 35, 40, 45\}$ entries (5 per count), logistic parameters drawn from GLUE/SuperGLUE ranges, Gaussian noise ($\sigma = 0.02$), random seed 42.

\subsection{Evaluation Metrics}

$R^2$ (fit quality), $\Delta$AIC (model preference), Spearman $\rho$ (temporal structure), saturation date error in months (application accuracy), convergence rate and mean $R^2$ at low $N$ (generalization).
